\documentclass[11pt]{letter}
\usepackage[margin=1in]{geometry}
\usepackage{hyperref}

\begin{document}

\begin{letter}{
Editor-in-Chief \\
IEEE Transactions on Dependable and Secure Computing \\
}

\opening{Dear Editor,}

We are pleased to submit our manuscript entitled \textbf{``Mind the Label Gap: How Fine-Tuned LLMs Hallucinate Sub-Categories in SOC Alert Classification''} for consideration in \emph{IEEE Transactions on Dependable and Secure Computing}.

\textbf{Summary.} We present the first systematic study of label compliance in fine-tuned LLMs for Security Operations Center (SOC) alert classification. While most models achieve perfect \emph{semantic} accuracy (normalized F1 = 100\%), we discover that strict label compliance varies from 74.9\% to 100\% depending on model architecture---a 25-point gap hidden by conventional evaluation metrics. Models hallucinate MITRE ATT\&CK sub-category names from pre-training knowledge, producing up to 4,947 schema-violating labels in a single evaluation. We term this phenomenon the ``label gap'' and introduce a dual-metric framework (strict vs.\ normalized F1) to expose it.

\textbf{Key contributions:}
\begin{enumerate}
\item A 7-model benchmark (0.8B--9B parameters) with strict and normalized F1 on clean, zero-overlap data splits.
\item A taxonomy of 10,093 label hallucination errors across 6 types (sub-technique substitution, morphological, cross-category, corrupted token, sibling class).
\item Discovery of non-monotonic scaling: 5K training samples produce 96$\times$ more errors than 1K, revealing a ``knowledge activation valley.''
\item Evidence that reasoning-oriented architectures (DeepSeek-R1, Phi4-mini) eliminate hallucination entirely, while model size does not predict compliance.
\item Cost analysis showing fine-tuning is 149--181$\times$ cheaper than commercial LLM APIs (GPT-4o, Claude).
\end{enumerate}

\textbf{Relevance to IEEE TDSC.} This work directly addresses dependability and security concerns in AI-assisted SOC operations. Label compliance failures in automated alert classification can trigger incorrect SOAR playbooks, distort threat landscape reporting, and compromise incident response. Our dual-metric framework provides a practical tool for evaluating the trustworthiness of fine-tuned LLMs in security-critical applications.

\textbf{Novelty.} While LLM hallucination has been extensively studied for factual generation, hallucination in \emph{classification tasks}---where models produce semantically correct but schema-violating labels---has received virtually no attention. We are the first to identify, categorize, and propose mitigations for this novel failure mode.

This manuscript has not been published previously and is not under consideration elsewhere. All authors have approved the manuscript and agree with its submission to IEEE TDSC.

\closing{Respectfully submitted,\\[2em]
Nutthakorn Chalaemwongwan\\
Department of Computer Engineering, KOSEN-KMITL\\
King Mongkut's Institute of Technology Ladkrabang\\
nutthakorn.ch@kmitl.ac.th
}

\end{letter}
\end{document}
